\documentclass[11pt,a4paper]{amsart}
\usepackage{amsmath,amssymb,amsthm,mathtools,hyperref}
\usepackage{geometry}
\geometry{margin=1in}

\theoremstyle{plain}
\newtheorem{theorem}{Theorem}[section]
\newtheorem{lemma}[theorem]{Lemma}
\newtheorem{proposition}[theorem]{Proposition}
\newtheorem{corollary}[theorem]{Corollary}

\theoremstyle{definition}
\newtheorem{definition}[theorem]{Definition}
\newtheorem{remark}[theorem]{Remark}
\newtheorem{example}[theorem]{Example}

\title{Post-Softmax Rank Expansion:\\A Counterexample in Transformer Analysis}
\author{Walker J. Kirkpatrick}
\date{June 20, 2026}
\address{Independent Research}
\email{drwjkirkpatrick-web@github}

\begin{document}
\maketitle

\begin{abstract}
It is commonly assumed in transformer analysis that the softmax operation
preserves or reduces the rank of a matrix. We show this assumption is
false by constructing an explicit family of rank-$1$ matrices whose
row-wise softmax has full rank. The simplest counterexample is a
$2\times 2$ matrix $Z = \begin{psmallmatrix}0&0\\0&ab\end{psmallmatrix}$
with rank $1$, whose softmax has rank $2$. We generalize to
$m\times n$ dimensions via a Vandermonde construction and discuss the
implications for low-rank attention approximations.
\end{abstract}

\section{Introduction}

Transformer architectures rely on the attention mechanism
\cite{vaswani2017attention}:
\[
\operatorname{Attention}(Q,K,V) = \operatorname{softmax}\!
\Bigl(\frac{QK^\top}{\sqrt{d_k}}\Bigr)V .
\]
Researchers frequently analyze the rank of pre-softmax and post-softmax
matrices. An implicit assumption is that the softmax, being an
element-wise monotone map, does not increase rank. We show this is
false.

\textbf{Main result.} There exist rank-$1$ matrices whose softmax has
rank~$\geq 2$ (and up to $\min(m,n)$ for suitably constructed
$m\times n$ matrices).

\section{Definitions}

For a real matrix $Z\in\mathbb{R}^{m\times n}$, the
\emph{row-wise softmax} is:
\[
\sigma(Z)_{ij} = \frac{\exp(Z_{ij})}{\sum_{k=1}^n \exp(Z_{ik})} .
\]
The \emph{numerical rank} with tolerance $\tau>0$ is:
\[
\operatorname{rank}_\tau(M) = \#
\{\,i\mid \sigma_i(M)\geq\tau\,\},
\]
where $\sigma_i(M)$ are the singular values of $M$ in descending order.

\section{The 2\texorpdfstring{$\times$}{x}2 Counterexample}

\begin{theorem}\label{thm:2x2}
There exists $Z\in\mathbb{R}^{2\times 2}$ with $\operatorname{rank}(Z)=1$
and $\operatorname{rank}(\sigma(Z))=2$.
\end{theorem}

\begin{proof}
Take
\begin{equation}\label{eq:Z2x2}
Z = \begin{pmatrix} 0 & 0 \\ 0 & ab \end{pmatrix}
= \begin{pmatrix} 0 \\ a \end{pmatrix}
\begin{pmatrix} 0 & b \end{pmatrix}
\end{equation}
with $a,b\neq 0$. Clearly $\operatorname{rank}(Z)=1$.

Row-wise softmax gives:
\[
\sigma(Z) = \begin{pmatrix}
\tfrac12 & \tfrac12 \\[4pt]
\tfrac{1}{1+e^{ab}} & \tfrac{e^{ab}}{1+e^{ab}}
\end{pmatrix} .
\]
Suppose $\alpha\cdot[\tfrac12,\tfrac12]+\beta\cdot
[\tfrac{1}{1+e^{ab}},\tfrac{e^{ab}}{1+e^{ab}}]=[0,0]$.
From the first component:
\[
\frac{\alpha}{2}+\frac{\beta}{1+e^{ab}}=0
\quad\Longrightarrow\quad
\beta = -\frac{\alpha(1+e^{ab})}{2}.
\]
Substituting into the second component:
\[
\frac{\alpha}{2}-\frac{\alpha e^{ab}}{2}=0
\quad\Longrightarrow\quad
\alpha(1-e^{ab})=0 .
\]
Since $ab\neq 0$, $e^{ab}\neq 1$, so $\alpha=0$ and $\beta=0$.
Thus the rows are linearly independent, proving
$\operatorname{rank}(\sigma(Z))=2$.
\end{proof}

\begin{example}
For $a=2$, $b=3$, we have $ab=6$ and
\[
\sigma(Z) = \begin{pmatrix}
0.500000 & 0.500000 \\
0.002473 & 0.997527
\end{pmatrix}.
\]
Numerical SVD gives singular values
$\sigma_1\approx 1.1426$, $\sigma_2\approx 0.4354$,
confirming rank~$2$ to machine precision.
\end{example}

\section{General \texorpdfstring{$m\times n$}{m x n} Construction}

\begin{theorem}\label{thm:general}
For any $m\geq 2$ and $n\geq 2$, there exists $Z\in\mathbb{R}^{m\times n}$
with $\operatorname{rank}(Z)=1$ and
$\operatorname{rank}(\sigma(Z))\geq 2$.
In particular, for $m,n\leq 5$, a Vandermonde construction gives
$\operatorname{rank}(\sigma(Z))=\min(m,n)$.
\end{theorem}

\begin{proof}
Let $u_i=a\cdot i$ and $v_j=b\cdot j$ for $i=0,\dots,m-1$,
$j=0,\dots,n-1$ with $a,b\neq 0$, and set $Z=uv^\top$. Then
\begin{align*}
\sigma(Z)_{ij}
&= \frac{\exp(ab\cdot ij)}{\sum_{k=0}^{n-1}\exp(ab\cdot ik)} \\
&= \frac{1}{S_i}\bigl(r^i\bigr)^j,
\end{align*}
where $r=e^{ab}>1$ and $S_i=\sum_{k=0}^{n-1}(r^i)^k$.

Each row $i$ is proportional to the geometric progression
$[1,r^i,r^{2i},\dots,r^{(n-1)i}]$. For $m\leq n\leq 5$,
these rows form a generalized Vandermonde matrix with distinct
nodes $r^i$. The determinant of any $m\times m$ submatrix is
$\prod_{0\leq i<j<m}(r^j-r^i)\neq 0$, so the rows are linearly
independent. Hence $\operatorname{rank}(\sigma(Z))=m=\min(m,n)$.
\end{proof}

\begin{remark}[Numerical Limitations]
For $m,n>5$, the exponential values $r^{ij}$ span ranges that exceed
double-precision arithmetic. The analytical result still holds in
exact arithmetic, but numerical verification requires arbitrary
precision or a modified construction.
\end{remark}

\section{Implications}

\textbf{Low-rank attention approximations.} Methods such as Performer
\cite{choromanski2021rethinking} and Linformer
\cite{wang2020linformer} approximate
$\sigma(QK^\top)\approx QK^\top$ by treating the softmax as
approximately linear. Our theorem shows this approximation has a
\emph{fundamental representational gap}: it cannot reproduce rank
expansions that the softmax creates.

\textbf{Practical consequence.} If attention is compressed to rank
$r$ \emph{before} the softmax, the softmax of the reconstruction may
still have higher rank than $r$. Correct compression must happen
\emph{after} the softmax.

\begin{corollary}
Any theorem assuming
$\operatorname{rank}(\sigma(Z))\leq\operatorname{rank}(Z)$ is
incorrect in general.
\end{corollary}

\section{Open Questions}

\begin{enumerate}
\item \textbf{Exact rank formula.} For $Z=uv^\top$, can the rank of
$\sigma(Z)$ be expressed in terms of the multiplicities of
components in $u$ and $v$?
\item \textbf{Approximate rank.} How close can $\sigma(Z)$ be to a
rank-$r$ matrix when $Z$ has rank $r$? Is there an Eckart--Young
bound for the softmax map?
\item \textbf{Arbitrary precision.} Can the Vandermonde construction
be verified for all $m,n$ using exact rational or symbolic arithmetic?
\end{enumerate}

\section*{Acknowledgments}
Verified empirically on NumPy with double-precision floating point.
CPU-only; no GPU required.

\begin{thebibliography}{9}
\bibitem{vaswani2017attention}
A. Vaswani et al., ``Attention is all you need,'' NeurIPS, 2017.

\bibitem{choromanski2021rethinking}
K. Choromanski et al., ``Rethinking attention with performers,''
ICLR, 2021.

\bibitem{wang2020linformer}
S. Wang et al., ``Linformer: Self-attention with linear complexity,''
arXiv:2006.04768, 2020.
\end{thebibliography}

\end{document}
